\section{Элементы теории случайных процессов}
\label{sec:random-processes}

\subsection{Определение случайного процесса и реализации}

\emph{Случайным процессом} $X(t)$ называется семейство случайных величин,
зависящих от неслучайного параметра $t$, который обычно интерпретируется как
время. При каждом фиксированном $t$ величина $X(t)$ является случайной
величиной. Если параметр принимает дискретные значения, получаем случайную
последовательность $X_1,X_2,\ldots$.

После проведения одного опыта зависимость $x(t)$ становится обычной
детерминированной функцией; она называется \emph{реализацией} случайного
процесса. Поэтому процесс можно рассматривать либо как семейство случайных
величин $X(t)$, либо как совокупность его возможных реализаций.

Полное вероятностное описание процесса задаётся всеми конечномерными законами
распределения
\[
    F(x_1,t_1;\ldots;x_n,t_n)
    =P\{X(t_1)<x_1,\ldots,X(t_n)<x_n\}
\]
для любых $n$ и любых наборов моментов $t_1,\ldots,t_n$. На практике такое
описание слишком громоздко, поэтому в корреляционной теории ограничиваются
моментами первого и второго порядков; при этом предполагается
$M\{|X(t)|^2\}<\infty$ для рассматриваемых моментов времени.

\subsection{Математическое ожидание, дисперсия и корреляционная функция}

Средний уровень процесса характеризуется функцией математического ожидания
\[
    m_X(t)=M\{X(t)\}.
\]
Дисперсия
\[
    D_X(t)=M\{[X(t)-m_X(t)]^2\}
\]
характеризует разброс значений процесса в данный момент времени.

Связь значений процесса в два момента времени описывает корреляционная
(автокорреляционная) функция
\[
    R_{XX}(t_1,t_2)
    =M\{[X(t_1)-m_X(t_1)][X(t_2)-m_X(t_2)]\}.
\]
Таким образом, $R_{XX}(t_1,t_2)$ --- это ковариация случайных величин
$X(t_1)$ и $X(t_2)$. В частности,
\[
    R_{XX}(t,t)=D_X(t).
\]
Она характеризует не только амплитуду разброса, но и степень временной связи,
то есть в определённой мере плавность случайного процесса.

\subsection{Стационарные случайные процессы}

Процесс называется \emph{стационарным в строгом смысле}, если его конечномерные
распределения не меняются при общем сдвиге времени:
\[
    F(x_1,t_1;\ldots;x_n,t_n)
    =F(x_1,t_1+t_0;\ldots;x_n,t_n+t_0)
\]
для любого $t_0$.

В корреляционной теории особенно важна \emph{стационарность в широком смысле}:
\[
    m_X(t)=m_X=\mathrm{const},
    \qquad
    R_{XX}(t_1,t_2)=R_{XX}(t_2-t_1)=R_{XX}(\tau).
\]
Тогда дисперсия постоянна:
\[
    D_X=R_{XX}(0),
\]
а корреляционная функция зависит только от временного сдвига. Для вещественного
стационарного процесса она чётна:
\[
    R_{XX}(\tau)=R_{XX}(-\tau),
\]
и по модулю не превосходит $R_{XX}(0)$.

\subsection{Сходимость и непрерывность в среднем квадратичном}

Для случайных величин и процессов естественна среднеквадратичная сходимость.
Говорят, что $X_n$ сходится к $X$ в среднем квадратичном, если
\[
    M\bigl[(X_n-X)^2\bigr]\to0.
\]

Процесс $X(t)$ называется \emph{непрерывным в среднем квадратичном} в точке
$t$, если
\[
    M\bigl[(X(t+\Delta)-X(t))^2\bigr]\to0,
    \qquad \Delta\to0.
\]
Раскрывая квадрат, получаем критерий через моментные характеристики:
\[
\begin{aligned}
M[(X(t+\Delta)-X(t))^2]
={}&D_X(t+\Delta)+D_X(t)\\
 &+[m_X(t+\Delta)-m_X(t)]^2
 -2R_{XX}(t,t+\Delta).
\end{aligned}
\]
Для стационарного процесса с постоянным средним это условие сводится к
непрерывности $R_{XX}(\tau)$ при $\tau=0$.

\subsection{Дифференцирование случайного процесса}

Производную также определяют в среднем квадратичном. Случайная величина
$\dot X(t)$ называется среднеквадратичной производной, если
\[
    \frac{X(t+\Delta)-X(t)}{\Delta}
    \longrightarrow \dot X(t)
\]
в среднем квадратичном при $\Delta\to0$. Такое определение не требует
дифференцируемости каждой отдельной реализации.

Существование производной связано с гладкостью корреляционной функции. В
частности, для стационарного процесса возможность среднеквадратичного
дифференцирования определяется поведением второй производной
$R_{XX}(\tau)$ в нуле; при существовании производной
\[
    D_{\dot X}=-R_{XX}''(0).
\]

\subsection{Гауссовские и марковские процессы}

Процесс называется \emph{гауссовским}, если для любых моментов
$t_1,\ldots,t_n$ случайный вектор
\[
    (X(t_1),\ldots,X(t_n))
\]
имеет многомерное нормальное распределение. Для гауссовского процесса все
конечномерные распределения определяются математическим ожиданием и
корреляционной функцией. Поэтому для него корреляционное описание особенно
информативно.

Марковское свойство выражает отсутствие дополнительной информации о будущем в
далёком прошлом при известном настоящем: условное распределение будущего при
заданном текущем состоянии не зависит от более ранних состояний. Простейший
дискретный пример --- авторегрессионная последовательность
\[
    X_{k+1}=aX_k+\xi_k,
\]
где новые возмущения $\xi_k$ независимы от прошлых состояний.

\subsection{Спектральная плотность стационарного процесса}

Для стационарного процесса временную структуру можно описывать не только
корреляционной функцией, но и в частотной области. В общем случае теорема
Винера--Хинчина связывает корреляционную функцию со спектральной мерой.
Если, в частности, $R_{XX}\in L^1(\mathbb R)$, спектральная мера имеет обычную
плотность, которую можно записать как преобразование Фурье:
\[
    S_{XX}(\omega)
    =\int_{-\infty}^{\infty}R_{XX}(\tau)e^{-i\omega\tau}\,d\tau.
\]
При выполнении условий существования обратного преобразования имеем
\[
    R_{XX}(\tau)
    =\frac{1}{2\pi}\int_{-\infty}^{\infty}
      S_{XX}(\omega)e^{i\omega\tau}\,d\omega.
\]
В частности,
\[
    D_X=R_{XX}(0)
    =\frac{1}{2\pi}\int_{-\infty}^{\infty}S_{XX}(\omega)\,d\omega.
\]
Для вещественного стационарного процесса $S_{XX}(\omega)$ вещественна,
неотрицательна и чётна. Формально постоянной спектральной плотности соответствует
идеальный белый шум; в непрерывном времени это обобщённый случай (его
корреляционная функция пропорциональна дельта-функции), поэтому идеальный белый
шум следует понимать как математическую идеализацию.

\subsection{Что важно сказать на экзамене}

Следует определить случайный процесс как семейство случайных величин и пояснить
понятие реализации. Полное описание --- все конечномерные распределения, но в
корреляционной теории используют $m_X(t)$, $D_X(t)$ и
$R_{XX}(t_1,t_2)$. Для стационарного в широком смысле процесса нужно обязательно
записать $m_X=\mathrm{const}$ и $R_{XX}=R_{XX}(\tau)$. Затем достаточно
объяснить среднеквадратичную непрерывность и дифференцируемость, привести
определения гауссовского и марковского процессов и закончить связью
корреляционной функции со спектральной плотностью через преобразование Фурье.

\newpage
